\documentclass[10pt,a4paper]{amsart}
\usepackage[margin=19mm]{geometry}
\usepackage{amsmath,amssymb,mathtools}
\usepackage[scr=rsfs,cal=euler]{mathalfa}
\usepackage{xfrac}
\usepackage[unicode,hidelinks,bookmarks=false]{hyperref}
\newcommand{\ie}{i.e.\ }
\newcommand{\cf}{cf.\ }
\newcommand{\resp}{respectively\ }
\newcommand{\Subsection}{\subsection}
\providecommand{\nobreakdash}{\nobreak\hbox{-}}
\setlength{\emergencystretch}{2em}
\multlinegap0pt
\usepackage{amssymb}
\usepackage[shortlabels]{enumitem}
\newtheorem{proposition}{Proposition}
\title{Some developments of exchangeable measure-valued P\'{o}lya sequences}
\author{Yoana R. Chorbadzhiyska, Hristo Sariev, Mladen Savov}
\date{Source: arXiv:2505.01594v3 (CC BY 4.0)}
\begin{document}
\maketitle

\noindent\emph{Source and licence.} Some developments of exchangeable measure-valued P\'{o}lya sequences. Version arXiv:2505.01594v3 (CC BY 4.0), distributed by arXiv under the Creative Commons Attribution 4.0 International licence, http://creativecommons.org/licenses/by/4.0/, as recorded in the arXiv licence metadata for this exact version and shown on the abstract page https://arxiv.org/abs/2505.01594v3. Full licence notice: \texttt{licenses/arxiv-2505.01594-CC-BY-4.0.txt}.\par\smallskip

\noindent\textit{Editorial scope.} Counted unit: the article's proposition results:other:cid (source0.tex lines 606--608). The article's complete original proof of it is delivered with it (source0.tex lines 1025--1079, one \texttt{proof} environment of 55 lines, which the article places away from the statement). No mathematics is written, repaired or re-derived: the statement and the proof are the article's own lines, in the article's own notation, and no retained line is substituted or re-flowed.  Retained uncounted as the source-local context the proof needs: lines 134--134; lines 171--173; lines 175--177; lines 350--352; lines 433--435.  Nothing was omitted from any retained span, so the package carries no omission record.  No retained line contains a citation, so the wrapper carries no bibliography and no bibliography entry.  No retained line contains a figure environment, an image-inclusion command or drawing code, so no image is delivered and the package declares no asset dependency.  Editorial formatting only, in addition: an AMS \texttt{amsart} preamble that restates the article's own declarations and macros used by the retained lines, namely the environments \texttt{proposition} on separate counters, together with the standard packages those lines need.  The retained environments print in this document as Proposition 1 in order of appearance, whereas the article prints the same items with the numbering of its own full document.  The complete native source of the article as distributed in the arXiv e-print for this exact version is bundled unchanged as \texttt{sources/177514/source0.tex}.
\subsection{Atoms of $\sigma$-algebras}\label{section:model:atoms}

\begin{equation}\label{condition:stationarity}\tag{$B$}
\int_\mathbb{X}R_x(A)\nu(dx)=\nu(A),\qquad\mbox{for all }A\in\mathcal{X};
\end{equation}

\begin{equation}\label{condition:reversible}\tag{$C$}
\int_\mathbb{X}R_y(A)R_x(dy)=R_x(A),\qquad\mbox{for all }A\in\mathcal{X}\mbox{ and }\nu\mbox{-a.e. }x.
\end{equation}

\begin{equation}\label{introduction_thrm}
\frac{R_x(\cdot)}{R_x(\mathbb{X})}=\nu(\cdot\mid\mathcal{G})(x)\qquad\mbox{for }\nu\mbox{-a.e. }x.
\end{equation}

\begin{proposition}\label{atoms_general}
Let $(X_n)_{n\geq 1}$ be an exchangeable MVPS*$(\theta,\nu,R)$ with strictly positive reinforcement. Take $\mathcal{G}$ to be the sub-$\sigma$-algebra in \eqref{introduction_thrm}, and define $\pi$ as in Section \ref{section:model:atoms} w.r.t. $\mathcal{G}$. Then $(\pi(X_n))_{n\geq 1}$ is a PS.
\end{proposition}

\begin{proposition}\label{results:other:cid}
Let $(X_n)_{n\geq1}$ be a balanced MVPS$(\theta,\nu,R)$. Then $(X_n)_{n\geq1}$ is c.i.d. if and only if $R$ satisfies \eqref{condition:stationarity} and \eqref{condition:reversible}.
\end{proposition}

\begin{proof}[Proof of Proposition \ref{results:other:cid}]
Suppose that $(X_n)_{n\geq1}$ is c.i.d. Let $A,B\in\mathcal{X}$. Then
\begin{equation}\label{eq:cid:proof:eq1}
\nu(A)=\mathbb{P}(X_1\in A)=\mathbb{P}(X_2\in A)=\int_\mathbb{X}\frac{\theta\nu(A)+R_x(A)}{\theta+1}\nu(dx),
\end{equation}
which after some simple algebra becomes
\begin{equation}\label{eq:cid:equality1-proof}
\int_\mathbb{X}R_x(A)\nu(dx)=\nu(A).
\end{equation}
On the other hand, we have
\begingroup\allowdisplaybreaks
\begin{align*}
\mathbb{P}(X_1\in A,X_2\in B)&=\int_A\frac{\theta\nu(B)+R_x(B)}{\theta+1}\nu(dx)=\frac{\theta\nu(A)\nu(B)+\int_AR_x(B)\nu(dx)}{\theta+1},
\end{align*}
\endgroup
and
\begingroup\allowdisplaybreaks
\begin{align*}
\mathbb{P}(X_1\in A,X_3\in B)&=\int_A\int_\mathbb{X}\frac{\theta\nu(B)+R_x(B)+R_y(B)}{\theta+2}\frac{\theta\nu(dy)+R_x(dy)}{\theta+1}\nu(dx)\\
&\begin{aligned}\;=\frac{1}{(\theta+2)(\theta+1)}&\biggl\{\theta^2\nu(A)\nu(B)+\theta\int_AR_x(B)\nu(dx)+\theta\nu(A)\int_\mathbb{X}R_y(B)\nu(dy)\\
&+\theta\nu(A)\nu(B)+\int_AR_x(B)\nu(dx)+\int_A\Bigl(\int_\mathbb{X}R_y(B)R_x(dy)\Bigr)\nu(dx)\biggr\}.\end{aligned}
\end{align*}
\endgroup
Since $(X_1,X_2)\overset{d}{=}(X_1,X_3)$, using \eqref{eq:cid:equality1-proof}, we obtain
\begingroup\allowdisplaybreaks
\begin{align*}
(\theta+2)&\biggl(\theta\nu(A)\nu(B)+\int_AR_x(B)\nu(dx)\biggr)\\
&=\theta^2\nu(A)\nu(B)+\theta\int_AR_x(B)\nu(dx)+2\theta\nu(A)\nu(B)+\int_AR_x(B)\nu(dx)+\int_A\Bigl(\int_\mathbb{X}R_y(B)R_x(dy)\Bigr)\nu(dx),
\end{align*}
\endgroup
or, after simplification, $\int_AR_x(B)\nu(dx)=\int_A(\int_\mathbb{X}R_y(B)R_x(dy))\nu(dx)$,
which implies that
\[R_x(B)=\int_\mathbb{X}R_y(B)R_x(dy),\qquad\mbox{for }\nu\mbox{-a.e. }x.\]\\

Conversely, suppose that $R$ satisfies \eqref{condition:stationarity} and \eqref{condition:reversible}. Repeating the argument in \eqref{eq:cid:proof:eq1} in reverse order, we get $X_1\overset{d}{=}X_2$. Moreover, for every $A\in\mathcal{X}$,
\begingroup\allowdisplaybreaks
\begin{align*}
\mathbb{P}(X_3\in A)&=\int_{\mathbb{X}^2}\frac{\theta\nu(A)+R_{x_1}(A)+R_{x_2}(A)}{\theta+2}\mathbb{P}(X_1\in dx_1,X_2\in dx_2)\\
&=\frac{1}{\theta+2}\biggl(\theta\nu(A)+\int_\mathbb{X}R_{x_1}(A)\nu(dx_1)+\int_\mathbb{X}R_{x_2}(A)\nu(dx_2)\biggr)=\nu(A);
\end{align*}
\endgroup
therefore, by induction, $(X_n)_{n\geq1}$ are identically distributed with marginal distribution $\nu$.

Fix $n\in\mathbb{N}$ and $A\in\mathcal{X}$. Let $C\in\mathcal{X}$ with $\nu(C)=1$ be the essential set in \eqref{condition:reversible}. Since $(X_n)_{n\geq1}$ are identically distributed with marginal distribution $\nu$, we have $\mathbb{P}(X_1\in C,\ldots,X_n\in C)=1$, see the proof of Proposition \ref{atoms_general}. It follows from \eqref{condition:stationarity} and \eqref{condition:reversible} that, for $\mathbb{P}_{(X_1,\ldots,X_n)}$-a.e. $(x_1,\ldots,x_n)\in\mathbb{X}^n$,
\begingroup\allowdisplaybreaks
\begin{align*}
\mathbb{P}(X_{n+2}\in A|X_1=x_1,&\ldots,X_n=x_n)=\int_\mathbb{X}\frac{\theta\nu(A)+\sum_{i=1}^{n+1}R_{x_i}(A)}{\theta+n+1}\mathbb{P}(X_{n+1}\in dx_{n+1}|X_1=x_1,\ldots,X_n=x_n)\\
&=\frac{1}{\theta+n+1}\bigg(\theta\nu(A)+\sum_{i=1}^nR_{x_i}(A)+\int_\mathbb{X}R_{x_{n+1}}(A)\frac{\theta\nu(dx_{n+1})+\sum_{i=1}^nR_{x_i}(dx_{n+1})}{\theta+n}\biggr)\\
&=\frac{1}{\theta+n+1}\bigg(\theta\nu(A)+\sum_{i=1}^nR_{x_i}(A)+\frac{\theta\nu(A)+\sum_{i=1}^nR_{x_i}(A)}{\theta+n}\biggr)\\
&=\frac{\theta\nu(A)+\sum_{i=1}^nR_{x_i}(A)}{\theta+n}\\
&=\mathbb{P}(X_{n+1}\in A|X_1,\ldots,X_n),
\end{align*}
\endgroup
which concludes the proof.
\end{proof}

\end{document}
